% !TeX program = lualatex
\documentclass[12pt,a4paper]{article}

% ========= حزمة =========
\usepackage[english, bidi=default, layout=counters]{babel}
\usepackage{amsmath,amssymb,amsfonts,mathtools}
\usepackage{booktabs}
\usepackage{tabularx}
\usepackage{array}
\usepackage{hyperref}
\usepackage{url}
\usepackage{graphicx}
\usepackage{xcolor}
\usepackage{titlesec}
\usepackage{fancyhdr}
\usepackage{microtype}
\usepackage{fontspec}
\usepackage{parskip}
\usepackage{float}
\usepackage{physics}
\usepackage{bm}
\usepackage{geometry}
\geometry{left=1.5cm,right=1.5cm,top=2.5cm,bottom=2.5cm}

% ========= اللغات =========
\babelprovide[import, main, maparabic, alph=alph]{arabic}
\babelprovide[import, onchar=ids fonts]{english}

% ========= الخطوط العربية =========
\defaultfontfeatures{Ligatures=TeX}
\setmainfont{Amiri}[
Script=Arabic,
Mapping=arabicdigits,
Ligatures=TeX,
BoldFont=Amiri-Bold,
ItalicFont=Amiri-Italic,
BoldItalicFont=Amiri-BoldItalic
]
\setsansfont{Amiri}[
Script=Arabic,
Mapping=arabicdigits
]
\newfontfamily{\arabicfonttt}{Amiri}[
Script=Arabic,
Mapping=arabicdigits
]

% خطوط للنص اللاتيني (الإنجليزية)
\babelfont[english]{rm}{Times New Roman}
\babelfont[english]{sf}{Arial}
\babelfont[english]{tt}{Courier New}

% ========= ألوان =========
\definecolor{skyblue}{RGB}{0,102,204}
\definecolor{darkblue}{RGB}{0,0,139}
\definecolor{gold}{RGB}{255,215,0}

% ========= أوامر YUCT (تمت إزالة \Tr لأنه موجود في physics) =========
\DeclareMathOperator{\Ent}{Ent}
% \DeclareMathOperator{\Tr}{Tr} % <-- محذوف: يُعرِّفها physics
\newcommand{\Dir}{\mathcal{D}}
\newcommand{\cL}{\mathcal{L}}
\newcommand{\Planck}{\mathrm{Pl}}
\newcommand{\Keff}{K_{\mathrm{eff}}}
\newcommand{\kappac}{\kappa_c}
\newcommand{\eps}{\varepsilon}
\newcommand{\betaf}{\beta}
\newcommand{\df}{d_{\!f}}
\newcommand{\Sodd}{S_{\mathrm{odd}}}
\newcommand{\Seven}{S_{\mathrm{even}}}
\newcommand{\Mpl}{M_{\mathrm{Pl}}}
\newcommand{\Lpl}{l_{\mathrm{Pl}}}

% ========= تنسيق العناوين =========
\titleformat{\section}{\Large\bfseries}{\thesection}{1em}{}
\titleformat{\subsection}{\large\bfseries}{\thesubsection}{1em}{}

% ========= رؤوس وتذييلات =========
\fancypagestyle{firstpage}{%
	\fancyhf{}%
	\renewcommand{\headrulewidth}{0pt}%
	\renewcommand{\footrulewidth}{0pt}%
	\fancyfoot[C]{%
		\begin{minipage}[t]{\textwidth}
			\centering
			\textcolor{gold}{\rule{\textwidth}{1.6pt}}\\[5pt]
			{\small \textsuperscript{1}أبحاث ياكوشيف، \url{https://yuct.org/}} \hfill
			{\small بروتوكول YPSDC: \url{https://ypsdc.com/}}\\[-2pt]
			\textcolor{gold}{\rule{\textwidth}{1.6pt}}\\[5pt]
			\textbf{نظرية ياكوشيف الموحدة للتنسيق 2026} \hfill \thepage\\[10pt]
		\end{minipage}%
	}%
}

\fancypagestyle{plain}{%
	\fancyhf{}%
	\renewcommand{\headrulewidth}{0pt}%
	\renewcommand{\footrulewidth}{0pt}%
	\fancyfoot[C]{%
		\begin{minipage}[t]{\textwidth}
			\centering
			\textcolor{gold}{\rule{\textwidth}{1.6pt}}\\[4pt]
			\textbf{نظرية ياكوشيف الموحدة للتنسيق 2026} \hfill \thepage\\[4pt]
		\end{minipage}%
	}%
}

\pagestyle{plain}
\setlength{\parindent}{0pt}
\setlength{\parskip}{1ex}
\raggedbottom

\hypersetup{
	colorlinks=true,
	linkcolor=blue,
	citecolor=blue,
	urlcolor=blue,
	pdftitle={YUCT والبنى المبددة: منظور قائم على التنسيق},
	pdfauthor={أليكسي ف. ياكوشيف}
}

% ========= رأس المقال =========
\newcommand{\makeheader}{%
	\noindent
	\begin{minipage}{0.17\textwidth}
		\raggedright
		{\sffamily\bfseries\color{skyblue}\fontsize{40}{40}\selectfont YUCT}%
	\end{minipage}%
	\hspace{30pt}
	\begin{minipage}{0.32\textwidth}
		\raggedright
		{\sffamily\fontsize{11}{13}\selectfont نظرية ياكوشيف\\ الموحدة\\ للتنسيق}%
	\end{minipage}%
	\hfill
	\begin{minipage}{0.38\textwidth}
		\raggedleft
		{\sffamily\bfseries ملحق}\\
		{\sffamily مايو 2026}\\
		{\sffamily\footnotesize \scriptsize\url{https://doi.org/10.5281/zenodo.18444598}}%
	\end{minipage}
	\par\vspace{5pt}
	\noindent\textcolor{gold}{\rule{\textwidth}{1.6pt}}
	\par\vspace{0pt}
}

\begin{document}
	\thispagestyle{firstpage}
	\makeheader
	
	\vspace{0cm}
	{\LARGE\bfseries YUCT والبنى المبددة \\ 
		\Large منظور قائم على التنسيق لعدم استقرار اللاتوازن \\[0.3cm]
		\large \textit{ربط قانون الخطأ الشامل ببداية الحمل الحراري} \par}
	\vspace{0.2cm}
	
	{\large أليكسي ف. ياكوشيف\footnote{أبحاث ياكوشيف، \url{https://yuct.org/}}} \\
	\vspace{0.0cm}
	
	{\color{darkblue}\textbf{\LARGE المستخلص}} \par
	{\large
		\noindent
		تصف نظرية ياكوشيف الموحدة للتنسيق (YUCT) التعقيد المنظم من خلال معامل وحيد هو كفاءة التنسيق \(K_{\mathrm{eff}}\) وقانون قياس شامل للخطأ \(\varepsilon = \kappa_c\alpha K_{\mathrm{eff}}^{-\beta}\) (\(\beta = 2/3\)، \(\kappa_c = 1/3\)). يستكشف هذا الملحق نتائج هذا القانون على استقرار الحالات المستقرة غير المتوازنة، ولا سيما بداية حمل بينار. نقترح تعريفاً لـ \(K_{\mathrm{eff}}\) لمائع بسيط بدلالة ناقليته الحرارية ولزوجته القصية — وهي كميات قابلة للقياس روتينياً. باستخدام هذا التعريف، نعبر عن معدل إنتاج الإنتروبيا \(\sigma\) كدالة لـ \(K_{\mathrm{eff}}\) ونُظهر أن عتبة الاستقرار الخطي تقابل قيمة حرجة شاملة لكفاءة التنسيق، \(K_{\mathrm{eff}}^{\mathrm{crit}} = (S_{\mathrm{odd}}/S_{\mathrm{even}})^{1/\beta} \approx 1.837\). تتبع هذه القيمة مباشرة من الحلقة الجبرية لـ YUCT (\(S_{\mathrm{odd}} = 6/5,\; S_{\mathrm{even}} = 4/5,\; \beta = 2/3\)) ولا تحتوي على معاملات حرة. التنبؤ هو أنه بالنسبة لأي مائع نيوتوني، يبدأ الحمل عندما تنخفض كفاءة تنسيقه الفعالة، كما هي معرفة هنا، إلى أقل من $\approx 1.837$. نشتق علاقة كمية قابلة للتكذيب بين عدد رايلي الحرج وعدد برانتل للمائع، يمكن اختبارها ببيانات مختبرية موجودة. لا يستخدم الاشتقاق علاقة KPZ أو الأبعاد الكسيرية؛ بل يعتمد فقط على البنية الجبرية لـ YUCT. كما نناقش بوضوح حدود المعالجة الظواهرية الحالية.
	}
	
	\vspace{0.1cm}
	\noindent
	{\color{darkblue}\textbf{الكلمات المفتاحية:}} YUCT، حمل بينار، كفاءة التنسيق، إنتاج الإنتروبيا، عدد برانتل، الحلقة الجبرية.
	
	\clearpage
	\tableofcontents
	\clearpage
	
	\section{مقدمة}
	تستند نظرية ياكوشيف الموحدة للتنسيق (YUCT) إلى فرضية أن التنسيق — قدرة مكونات نظام ما على مواءمة حالاتها باستخدام قواميس مشتركة مسبقة — هو العملية الأساسية التي تنبثق منها القوانين الفيزيائية~\cite{Yakushev2026}. تكمن قوتها التنبؤية في قانون قياس وحيد للخطأ النسبي للتنسيق:
	\begin{equation}\label{eq:universal}
		\varepsilon = \kappa_c\,\alpha\, K_{\mathrm{eff}}^{-\beta},\qquad 
		\beta = \frac{2}{3},\quad \kappa_c = \frac{1}{3},
	\end{equation}
	حيث \(\alpha\) ثابت خاص بالمنظومة من مرتبة الوحدة. قيم \(\beta\) و \(\kappa_c\) ليست معاملات حرة؛ بل تتبع من الاتساق الذاتي الجبري لشبكات التنسيق الهرمية، كما هو مفصل في الملاحق~\cite{AppendixY, AppendixL, AppendixFerra, AppendixAF}. تنتج الحلقة الجبرية الناتجة عن ذلك ثابتين كسريين مضبوطين،
	\begin{equation}\label{eq:S-constants}
		S_{\mathrm{odd}} = \frac{6}{5} = 1.2,\qquad
		S_{\mathrm{even}} = \frac{4}{5} = 0.8,
	\end{equation}
	يحققان \(S_{\mathrm{odd}} + S_{\mathrm{even}} = 2\) و \(S_{\mathrm{odd}}/S_{\mathrm{even}} = 3/2 = 1/\beta\). هذه الأعداد هي العمود الفقري لكل الاشتقاقات اللاحقة.
	
	في خط بحثي مختلف، تصف نظرية البنى المبددة التي صاغها إيليا بريغوجين كيف يمكن للمنظومات المفتوحة المدفوعة بعيداً عن التوازن الترموديناميكي أن تتطور تلقائياً نحو حالات أعلى تعقيداً~\cite{Prigogine1977, Glansdorff1971}. المثال النموذجي هو حمل بينار: طبقة مائع أفقية مسخنة من الأسفل تُطور نمطاً منتظماً من لفات الحمل عندما يتجاوز الممال الحراري قيمة حرجة~\cite{Chandrasekhar1961}. ينص معيار بريغوجين للاستقرار على أن الحالة المستقرة تصبح غير مستقرة عندما يصبح فائض إنتاج الإنتروبيا \(\delta_X P\) موجباً.
	
	الغرض من هذا الملحق هو إظهار أن معيار بريغوجين يمكن إعادة صياغته بلغة YUCT، وأن الحلقة الجبرية لـ YUCT تقدم تنبؤاً كمياً قابلاً للاختبار لبداية الحمل. الاشتقاق ظواهري، لكنه لا يحتوي على معاملات اعتباطية تتجاوز تلك الموجودة أصلاً في إطار YUCT. كما نقر صراحة بحدود المقاربة الحالية ونوجز الخطوات اللازمة لنظرية مجهرية كاملة.
	
	\section{كفاءة تنسيق مائع بسيط}
	لتطبيق YUCT على منظومة فيزيائية، يجب تحديد «القاموس» و«الدليل» المناسبين، ومن ثم حساب \(K_{\mathrm{eff}}\). بالنسبة لمائع في ممال حراري، يمكن النظر إلى نقل الحرارة كسلسلة من أحداث تنسيق موضعية: تنقل الجزيئات عند درجة حرارة أعلى الطاقة إلى تلك عند درجة حرارة أخفض عبر التصادمات. كفاءة هذه العملية محدودة بعاملين: القدرة الذاتية للوسط على توصيل الحرارة (ناقليته الحرارية \(\lambda\)) والاحتكاك الداخلي الذي يبدد الطاقة (لزوجته القصية \(\eta\)). المائع ذو \(\lambda\) عالية و \(\eta\) منخفضة هو «منسق جيد» لأن الطاقة تُنقل بسرعة مع خسارة قليلة.
	
	لذلك نقترح التعريف العملياتي التالي:
	\begin{equation}\label{eq:Keff-fluid}
		K_{\mathrm{eff}} \equiv \frac{\lambda}{\lambda_0} \left( \frac{\eta_0}{\eta} \right)^{\gamma},
	\end{equation}
	حيث \(\lambda_0\) و \(\eta_0\) قيمتان مرجعيتان (مثلاً للماء عند $20^\circ$م) ويُختار الأُس \(\gamma\) بحيث يصبح \(K_{\mathrm{eff}}\) كمية لابعدية ومستقلة عن المقياس. بما أن النسبة \(\lambda/\eta\) لها بعد \((\mathrm{طول})^2/(\mathrm{زمن}\cdot\mathrm{درجة حرارة})\)، فإن الاختيار الطبيعي هو \(\gamma = 1\)، مما يعطي
	\begin{equation}\label{eq:Keff-fluid2}
		K_{\mathrm{eff}} = \frac{\lambda / \lambda_0}{\eta / \eta_0}.
	\end{equation}
	هذا التعريف مشابه لـ«عدد التنسيق» المقدم للعناصر الكيميائية في الملحق~C، ويجعل \(K_{\mathrm{eff}}\) قابلاً للقياس مباشرة.
	
	\section{إنتاج الإنتروبيا وقانون الخطأ الشامل}
	بالنسبة لمائع ساكن في ممال حراري \(\nabla T\)، فإن التدفق الترموديناميكي الوحيد هو تدفق الحرارة \(J = -\lambda \nabla T\)، والقوة المرافقة هي \(X = \nabla(1/T)\). معدل إنتاج الإنتروبيا لكل وحدة حجم هو
	\begin{equation}\label{eq:sigma}
		\sigma = J \cdot \nabla\left(\frac{1}{T}\right) \approx \frac{\lambda |\nabla T|^2}{T^2}.
	\end{equation}
	في صورة YUCT، تدفق الحرارة هو نتيجة لعدد كبير من أفعال التنسيق المجهرية. جزء \(\varepsilon\) من هذه الأفعال يفشل في نقل الطاقة بشكل مترابط، لذا ينخفض التدفق الفعال:
	\begin{equation}\label{eq:J-effective}
		J_{\mathrm{eff}} = J \; \bigl(1 - \varepsilon\bigr).
	\end{equation}
	يساهم التدفق «المفقود» في إنتاج الإنتروبيا ولكن ليس في الإزالة المترابطة للممال. باستخدام قانون الخطأ الشامل~\eqref{eq:universal}، يمكننا كتابة
	\begin{equation}\label{eq:sigma-K}
		\sigma = \frac{\lambda |\nabla T|^2}{T^2} \frac{1}{1 - \kappa_c\alpha K_{\mathrm{eff}}^{-\beta}}.
	\end{equation}
	عندما يكون \(K_{\mathrm{eff}} \gg 1\) يكون التصحيح مهملاً، وتطيع المنظومة قانون فورييه الخطي المعتاد. كلما انخفض \(K_{\mathrm{eff}}\) (أصبح المائع ألزج بالنسبة لناقليته)، ازداد الخطأ، وازداد إنتاج الإنتروبيا لنفس الممال المفروض.
	
	\section{عتبة الاستقرار وكفاءة التنسيق الحرجة}
	ينص معيار بريغوجين للاستقرار على أن الحالة المستقرة تصبح غير مستقرة عندما يتسبب تذبذب في جعل فائض إنتاج الإنتروبيا موجباً~\cite{Glansdorff1971}. في صياغة YUCT، يمكن التفكير في التذبذب كتغير موضعي في حالة تنسيق المائع — انحراف مؤقت لـ \(K_{\mathrm{eff}}\) عن قيمته المستقرة. باستخدام~\eqref{eq:sigma-K}، يؤدي تغير صغير \(\delta K_{\mathrm{eff}}\) إلى تغير في إنتاج الإنتروبيا:
	\begin{equation}
		\delta_X P \propto - \delta K_{\mathrm{eff}}.
	\end{equation}
	وهكذا، فإن التذبذب الذي \emph{يخفض} \(K_{\mathrm{eff}}\) يجعل \(\delta_X P > 0\) ويميل إلى أن يتضخم. هذا التضخيم يسرع تدهور التنسيق، دافعاً المنظومة نحو حالة جديدة.
	
	لذلك تصبح حالة التوصيل المستقرة غير مستقرة عندما ينخفض \(K_{\mathrm{eff}}\) تحت قيمة حرجة. لأن الحلقة الجبرية لـ YUCT تحدد جميع النسب اللابعدية، يجب أن تكون هذه القيمة الحرجة قابلة للتعبير عنها بدلالة الثوابت الأساسية:
	\begin{equation}\label{eq:Kcrit}
		K_{\mathrm{eff}}^{\mathrm{crit}} = \left( \frac{S_{\mathrm{odd}}}{S_{\mathrm{even}}} \right)^{1/\beta}
		= \left( \frac{3}{2} \right)^{3/2} \approx 1.837.
	\end{equation}
	يتبع الأُس \(1/\beta\) من شرط أن يكون معامل الخطأ \(\varepsilon\) عند العتبة من مرتبة الوحدة (حيث تعجز المنظومة عن تمييز التنسيق المفيد من الضوضاء). النسبة \(S_{\mathrm{odd}}/S_{\mathrm{even}} = 3/2\) هي نسبة النظام إلى الفوضى الشاملة في YUCT. لا تحتوي المعادلة~\eqref{eq:Kcrit} على أي معاملات حرة.
	
	\section{تنبؤ بعدد رايلي الحرج}
	عدد رايلي لطبقة مائع عمقها \(d\)، كثافتها \(\rho\)، حرارتها النوعية \(c_p\)، معامل تمددها الحراري \(\alpha_T\)، وفرق درجة الحرارة المفروض \(\Delta T\) هو
	\begin{equation}
		\mathrm{Ra} = \frac{\rho^2 c_p \alpha_T g \Delta T d^3}{\lambda \eta}.
	\end{equation}
	بدلالة كفاءة التنسيق لدينا،~\eqref{eq:Keff-fluid2}، يمكننا كتابة
	\begin{equation}
		\mathrm{Ra} = \frac{A}{K_{\mathrm{eff}}},
	\end{equation}
	حيث \(A\) ثابت يعتمد فقط على الهندسة والشروط الحدية، وليس على خواص المائع. بهذا التحديد، يترجم شرط العتبة \(K_{\mathrm{eff}} = K_{\mathrm{eff}}^{\mathrm{crit}}\) إلى عدد رايلي حرج:
	\begin{equation}\label{eq:Racrit}
		\mathrm{Ra_c} = \frac{A}{K_{\mathrm{eff}}^{\mathrm{crit}}}.
	\end{equation}
	إذا أمكن تحديد الثابت \(A\) من مائع مرجعي واحد، وليكن الماء، الذي له عدد رايلي الحرج التقليدي \(\mathrm{Ra}_{\mathrm{water}} \approx 1708\) (لأجل حدود صلبة) ويعطينا تعريفنا قيمة ما \(K_{\mathrm{eff}}^{\mathrm{water}}\)، عندئذ يمكن التنبؤ بعدد رايلي الحرج لأي مائع آخر:
	\begin{equation}\label{eq:Racrit-predict}
		\boxed{ \mathrm{Ra_c}(\text{مائع}) = \mathrm{Ra_c}(\text{ماء}) \;
			\frac{K_{\mathrm{eff}}^{\mathrm{water}}}{K_{\mathrm{eff}}(\text{مائع})} }.
	\end{equation}
	باستخدام~\eqref{eq:Keff-fluid2}، يصبح هذا علاقة بين عدد رايلي الحرج والناقلية الحرارية واللزوجة للمائع:
	\begin{equation}
		\mathrm{Ra_c} \propto \frac{\eta / \eta_0}{\lambda / \lambda_0}.
	\end{equation}
	وهكذا، ينبغي لمائع ذي لزوجة أعلى (تنسيق أضعف) أن يكون له عدد رايلي حرج \emph{أكبر} — يصعب جعله يحمل. هذا التنبؤ متسق نوعياً مع النتائج التقليدية: الزيوت، اللزجة، لها أعداد رايلي حرجة أعلى بكثير من الماء. يتطلب الاختبار الكمي قياس \(\lambda\) و \(\eta\) للمائع، وحساب \(K_{\mathrm{eff}}\) من~\eqref{eq:Keff-fluid2}، ومقارنة \(\mathrm{Ra_c}\) المتنبأ به مع بداية الحمل المرصودة تجريبياً. لا تبقى أي معاملات قابلة للتعديل بعد المعايرة الوحيدة للمائع المرجعي.
	
	\section{نتائج قابلة للاختبار وبروتوكول تجريبي}
	\subsection{اختبار فوري بالبيانات الموجودة}
	تحتوي المراجع على جداول موسعة لأعداد رايلي الحرجة لموائع مختلفة (انظر مثلاً المراجعات~\cite{Chandrasekhar1961, Cross1993}). باستخدام القيم المقيسة لـ \(\lambda\) و \(\eta\) عند درجات الحرارة المناسبة، يمكن حساب \(K_{\mathrm{eff}}\) لكل مائع والتحقق مما إذا كانت النسبة \(\mathrm{Ra_c}/\mathrm{Ra_c}^{\mathrm{water}}\) تساوي \(K_{\mathrm{eff}}^{\mathrm{water}} / K_{\mathrm{eff}}\). ارتباط موجب مع \(R^2 > 0.9\) سيقدم دليلاً قوياً على آلية YUCT.
	
	\subsection{تجربة مخصصة}
	يمكن إجراء تجربة مضبوطة باستخدام مائعين لهما كفاءتا تنسيق مختلفتان بشكل كبير، مثلاً الماء وزيت السيليكون. يُقاس فرق درجة الحرارة الحرج \(\Delta T_c\) لكل مائع بدقة عالية. تنبؤ YUCT هو أن نسبة فروق درجات الحرارة الحرجة (عند هندسة ثابتة) يجب أن تساوي النسبة العكسية لكفاءات التنسيق كما تعرفها~\eqref{eq:Keff-fluid2}. هذا تنبؤ خالٍ من المعاملات وقابل للتكذيب.
	
	\section{حدود وأسئلة مفتوحة}
	التحليل الحالي ظواهري. ستتطلب النظرية المجهرية الكاملة اشتقاق الصلة بين خطأ التنسيق \(\varepsilon\) والتدفقات الترموديناميكية من المبادئ الأولى، دون افتراض~\eqref{eq:J-effective}. علاوة على ذلك، تعريف \(K_{\mathrm{eff}}\) في~\eqref{eq:Keff-fluid2} هو حدس مبرر؛ قد تكون هناك تعريفات أخرى ممكنة وينبغي اختبارها مقابل البيانات. يتجاهل التحليل أيضاً دور التفاعلات الكيميائية، المعروف أنها تعدل بداية الحمل بقوة. هذه مواضيع للبحث المستقبلي.
	
	رغم هذه الحدود، فإن الاشتقاق الجبري لقيمة كفاءة تنسيق حرجة \(K_{\mathrm{eff}}^{\mathrm{crit}} \approx 1.837\) من حلقة YUCT هو مضبوط داخل الإطار. ما إذا كانت الطبيعة تستخدم هذه القيمة فعلاً يمكن أن يقرره التجريب.
	
	\section{استنتاجات}
	لقد أظهرنا أن إطار YUCT يقدم تفسيراً كمياً قابلاً للاختبار لبداية حمل بينار. يُعبر عن عدد رايلي الحرج بدلالة كفاءة تنسيق المائع، المرتبطة بدورها بمعاملات نقل قابلة للقياس. يمكن التحقق من علاقة القياس المتنبأ بها \(\mathrm{Ra_c} \propto \eta/\lambda\) باستخدام بيانات مختبرية موجودة. يوضح هذا العمل أن YUCT ليست مجرد إعادة تفسير فلسفية للترموديناميك، بل إطار قادر على توليد تنبؤات عددية قابلة للتكذيب.
	
	\section*{شكر وتقدير}
	يشكر المؤلف المشاركين في ندوة YUCT على المناقشات النقدية.
	
	\begin{thebibliography}{99}
		\bibitem{Yakushev2026} A.~V.~Yakushev, \emph{Yakushev Unified Coordination Theory (YUCT)}, Zenodo, DOI:10.5281/zenodo.18444599 (2026).
		\bibitem{AppendixY} A.~V.~Yakushev, ``Appendix Y: Mathematical Foundations of YUCT,'' in \emph{YUCT}, Zenodo (2026).
		\bibitem{AppendixL} A.~V.~Yakushev, ``Appendix L: Fractal Coordination Error Scaling,'' in \emph{YUCT}, Zenodo (2026).
		\bibitem{AppendixFerra} A.~V.~Yakushev, ``Appendix Ferra1.2: Universal Coordination Constants for Ordered and Disordered Phases,'' in \emph{YUCT}, Zenodo (2026).
		\bibitem{AppendixAF} A.~V.~Yakushev, ``Appendix AF: The Algebraic Framework of Reality in YUCT,'' in \emph{YUCT}, Zenodo (2026).
		\bibitem{Prigogine1977} I.~Prigogine and G.~Nicolis, \emph{Self‑Organization in Nonequilibrium Systems}. Wiley (1977).
		\bibitem{Glansdorff1971} P.~Glansdorff and I.~Prigogine, \emph{Thermodynamic Theory of Structure, Stability and Fluctuations}. Wiley (1971).
		\bibitem{Chandrasekhar1961} S.~Chandrasekhar, \emph{Hydrodynamic and Hydromagnetic Stability}. Oxford (1961).
		\bibitem{Cross1993} M.~C.~Cross and P.~C.~Hohenberg, ``Pattern formation outside of equilibrium,'' \emph{Rev. Mod. Phys.} \textbf{65}, 851 (1993).
	\end{thebibliography}
	
\end{document}